\documentclass[11pt]{report}
\input{../../../report.tex}
\title{Algebra Exercises}
\usepackage{titlepageBU}
\usepackage{xparse}
\courseID{CAS MA 741}
\college{College of Arts and Sciences}
\department{Department of Mathematics and Statistics}
\date{Winter Break, 2024/2025}
\usepackage{tikz-cd}
\DeclareMathOperator{\Stab}{Stab}
\DeclareMathOperator{\rank}{rk}
\DeclareMathOperator{\End}{End}
\DeclareMathOperator{\Hom}{Hom}
\DeclareMathOperator{\Obj}{Obj}
\DeclareMathOperator{\Aut}{Aut}
\DeclareMathOperator{\Tr}{tr}
\newcommand{\tildesim}{\mathbin{\sim}}
\declaretheoremstyle[headfont=\bfseries, bodyfont=\normalfont]{x}
\declaretheorem[numberwithin=section, style=x, name=Exercise]{exercise}
\declaretheorem[numbered=no, style=proof, name=Solution]{solution}
\newcommand{\contradiction}{\Longrightarrow\Longleftarrow}
\begin{document}
\makereport
\stepcounter{chapter}
\stepcounter{chapter}
\stepcounter{chapter}
\stepcounter{chapter}
\stepcounter{chapter}
\chapter{Linear Algebra}
\section{Free Modules Revisited}
\begin{exercise}
	Show $\mathbb{R} \cong \mathbb{C} $ as $\mathbb{Q} $-vector spaces.
\end{exercise}
\begin{solution}
	This is tricky. I think $\dim_{\mathbb{Q} } \mathbb{R} =\aleph_0$, and $\dim_{\mathbb{Q} } \mathbb{C} =2 \dim _{\mathbb{Q} } \mathbb{R} =\dim_{\mathbb{Q} } \mathbb{R} $. Unsure how to show this.
\end{solution}
\begin{exercise}
	Consider the sets
	\begin{itemize}
		\item $\mathfrak{gl}_n(\mathbb{R} )\coloneqq \mathcal{M}_n(\mathbb{R} )$;
	\item $\mathfrak{sl}_n(\mathbb{R} )\coloneqq \left\{ M\in \mathfrak{gl}_n(\mathbb{R} )\mid \Tr M=0 \right\} $;
	\item $\mathfrak{sl}_n(\mathbb{C} )\coloneqq \left\{ M\in \mathfrak{gl}_n(\mathbb{C} )\mid \Tr M=0 \right\} $;
	\item $\mathfrak{so}_n(\mathbb{R} )\coloneqq \left\{ M\in \mathfrak{gl}_n(\mathbb{R} )\mid M+M^t=0 \right\} $;
	\item $\mathfrak{su}(n)\coloneqq \left\{ M\in \mathfrak{gl}_n(\mathbb{C} )\mid M+M^\dagger=0 \right\} $.
	\end{itemize}
	Show these sets are $\mathbb{R} $-vector spaces, and compute their dimension.
\end{exercise}
\begin{solution}
	First, we show $\mathfrak{gl} _n(\mathbb{R} )$ is a $\mathbb{R} $-vector space. We will use the componentwise matrix notation $M\coloneqq [m_{ij} ]$, where $m_{ij} $ are the components of the matrix $M$ and range over $\left\{ 1, \ldots ,n \right\} $. For any $r\in\mathbb{R} $, define $r[m_{ij} ]=[r m_{ij} ] $. Clearly, $\mathfrak{gl} _n(\mathbb{R} )$ is a group, so we will simply show it is a vector space. Note that $1[m_{ij} ]=[m_{ij} ]$,
	\[
		(r+s)[m_{ij} ]=r[m_{ij} ]+s[m_{ij} ]
	,\]
	and
	\[
		(rs)[m_{ij} ]=[rsm_{ij} ]=r[sm_{ij} ]=r\left( s[m_{ij} ] \right)  
	.\]
	Moreover, $r$ is an endomorphism:
	\[
		r\left( [m_{ij} ]+[n_{ij} ] \right) =r[m_{ij} +n_{ij} ]=[r\left( m_{ij} +n_{ij}  \right) ]=[r m_{ij} +rn_{ij} ]=[r m_{ij} ]+[rn_{ij} ]=r[m_{ij} ]+r[n_{ij} ] 
	.\]
	Therefore, $r\mapsto r$ is a homomorphism $\mathbb{R}  \to \End_{\mathsf{Ab} }(\mathfrak{gl}_n(\mathbb{R} )) $, and $\mathfrak{gl} _n(\mathbb{R} )$ is an $\mathbb{R} $-vector space. Finally, we compute the dimension. Obviously, $\mathfrak{gl}_n(\mathbb{R} )\cong \mathbb{R} ^{n^2} $. This is so much effort to prove so I'm not going to.

	$\mathfrak{sl}_n(\mathbb{R} )$ is very difficult, but note that $\mathfrak{so}_n(\mathbb{R} )\subseteq \mathfrak{sl}_n(\mathbb{R} )\subseteq \mathfrak{gl}_n(\mathbb{R} )$,  so maybe that helps.

	I think the basis for $\mathfrak{so}_n(\mathbb{R} )$ is basically one element per item $m_{ij} $ in the upper triangle (not counting main diagonal), put a 1 there, and put a $-1$ at $m_{ji} $. Then $M^t=-M$ for the basis set. I think this is the entire basis.
\end{solution}

\end{document}

